% Part: lambda-calculus
% Chapter: representability
% Section: currying

\documentclass[../../../include/open-logic-section]{subfiles}

\begin{document}

\olfileid{lam}{rep}{cur}
\olsection{करीकरण}

$\lambd[x][M]$ हे $\lambd$-अमूर्तीकरण एकच फलसाधक घेणारे
फलन दर्शवते; अनेक फलसाधक घेणारी फलने परिभाषित करायची असतील,
तर ही मोठी मर्यादा ठरते. एक मार्ग म्हणजे जोड्या, तिकड्या इत्यादी
घडवता येतील असा $\lambd$-कलनाचा विस्तार करणे. मग, उदाहरणार्थ,
त्रिस्थानी फलन $\lambd[x][M]$ ला त्याचा फलसाधक म्हणून
तिकडी अपेक्षित असेल. मात्र हे \emph{करीकरणाने} करणे अधिक सोयीचे आहे.

एक उदाहरण पाहू. क्षणभर $\lambd$-कलनात $+$ ही क्रिया आहे
असे मानू. बेरीज हे $2$-स्थानी फलन आहे, म्हणजे ते दोन फलसाधक
घेते. पण $\lambd$-अमूर्तीकरण आपल्याला एकच फलसाधक घेणारी फलने
देते: $\lambd[(x,y)][(x+y)]$ यांसारख्या व्यक्तीकरणांना विन्यासात
परवानगी नाही. तरी $\lambd[y][(x+y)]$ ने दिलेले एकस्थानी
फलन~$f_x(y)$ घेता येते; ते आपल्या एकमेव फलसाधक~$y$ मध्ये
$x$ मिळवते. खरे तर हे एक फलन नाही, तर प्रत्येक संख्या~$x$
साठी “$x$ मिळवा” अशा वेगवेगळ्या फलनांचे कुटुंब आहे. आता
$\lambd[x][f_x]$ असे आणखी एक एकस्थानी फलन~$g$ परिभाषित
करता येते. फलसाधक $x$ ला उपयोजल्यास $g(x)$ हे
फलन~$f_x$ परत करते---म्हणजे $g$ ची मूल्ये इतर फलने आहेत.
आता $g$ चे $x$ ला आणि मिळालेल्या फलनाचे~$y$ ला उपयोजन
केल्यास $(g(x))y = f_x(y) = x+y$ मिळते. अशा प्रकारे हे एकस्थानी फलन द्विस्थानी बेरीज फलनाचेच काम करू शकते.
“करीकरण” म्हणजे द्विस्थानी फलनाचे, ज्याची मूल्ये एकस्थानी
फलने आहेत अशा, एकस्थानी फलनात रूपांतर करण्याची ही युक्ती.

%READERNOTE{SOL6-A105: या करीकरणाच्या न्यूनीकरणांत अमूर्त केलेली चरे परस्पर भिन्न आणि सर्व आदेशित पदांतील मुक्त चरांपासून वेगळी निवडली आहेत. आवश्यक नामांतर आधी केले जाते; अन्यथा displayed intermediate expression मुक्त चर पकडू शकते.}
आता $\lambd$-कलनाच्या विन्यासात बसणारे उदाहरण पाहू.
या उदाहरणातील अमूर्तीकरणाचे दुसरे चर आदेशित पहिल्या पदात
मुक्त नसावे; आवश्यक असल्यास त्याचे आधी नामांतर करा.
$f(x,y) = x$ हे फलन कसे दर्शवायचे? दोन फलसाधक घेऊन
पहिला परत करणारे फलन परिभाषित करायचे असेल, तर
$\lambd[x][\lambd[y][x]]$ लिहिता येते. हे अक्षरशः
फलसाधक~$x$ घेऊन फलन~$\lambd[y][x]$ परत करणारे फलन आहे.
$\lambd[y][x]$ हे फलन आणखी एक फलसाधक~$y$ घेते,
पण तो बाजूला ठेवून नेहमी~$x$ परत करते. दोन फलसाधकांना
$\lambd[x][\lambd[y][x]]$ उपयोजल्यास काय होते ते पाहू:
\begin{align*}
  (\lambd[x][\lambd[y][x]])MN
  \bredone & (\lambd[y][M])N \\
  \bredone & M
\end{align*}

पुढील सर्वसाधारण मांडणीत अमूर्तीकरणातील चरे परस्पर भिन्न
आणि उपयोजल्या जाणाऱ्या सर्व पदांतील मुक्त चरांपासून वेगळी
निवडली आहेत; आवश्यक नामांतर आधी केले आहे.
सर्वसाधारणपणे, $x_1$, \dots,~$x_n$ ही प्राचले असलेले
आणि पद~$N$ ने परिभाषित होणारे फलन लिहायचे असेल,
तर $\lambd[x_1][\lambd[x_2][\ldots\lambd[x_n][N]]]$
लिहिता येते. त्याला $n$ फलसाधक उपयोजल्यास मिळते:
\begin{multline*}
  (\lambd[x_1][\lambd[x_2][\ldots\lambd[x_n][N]]]) M_1 \dots M_n \bredone\\
  \begin{aligned}
  \bredone {} & (\Subst{(\lambd[x_2][\ldots\lambd[x_n][N]])}{M_1}{x_1}) M_2
  \dots M_n\\
   \eqs {} & (\lambd[x_2][\ldots\lambd[x_n][\Subst{N}{M_1}{x_1}]]) M_2
            \dots M_n \\
  \vdots & \\
  \bredone {} &\Subst{\Subst{N}{M_1}{x_1}\ldots}{M_n}{x_n}
  \end{aligned}
\end{multline*}
शेवटच्या ओळीचा अक्षरशः अर्थ फलनाच्या व्याख्येच्या देहात
$x_i$ च्या जागी $M_i$ आदेशणे असा आहे; फलनाला अनेक
फलसाधक उपयोजल्यावर आपल्याला नेमके हेच अपेक्षित असते.
\end{document}
